\documentclass[12pt]{article}
\usepackage[T1]{fontenc}
\usepackage[utf8]{inputenc}
\usepackage{lmodern}
\usepackage{textcomp}
\usepackage{enumitem}
\usepackage{geometry}
\usepackage{graphicx}
\usepackage{amsmath, amssymb}
\geometry{margin=1in}
\begin{document}
\begin{center}
History - Graduate Level Assessment 12 \
Total Marks: 40 \
Answer all questions.
\end{center}

\section*{Multiple Choice (10 marks)}
\begin{enumerate}[label=\arabic*.]
\item The native peoples of the northwest coast of North America were:
\begin{enumerate}[label=(\alph*)]
    \item simple foragers.
    \item maize agriculturalists.
    \item conquered by the Aztec.
    \item dependent on hunting large game.
    \item primarily seafaring traders.
\end{enumerate}
\item The peak of Minoan civilization, from 3650 to 3420 B.P., was bracketed by what two events?
\begin{enumerate}[label=(\alph*)]
    \item the introduction of Linear A script and the adoption of Linear B script
    \item an earthquake and a volcanic eruption
    \item the end of the Stone Age and the beginning of the Bronze Age
    \item the discovery of bronze on the island, and the replacement of bronze by iron
    \item the construction of the Palace of Knossos and its destruction
\end{enumerate}
\item According to one theory, the power of Ancestral Pueblo elites emerged when certain groups:
\begin{enumerate}[label=(\alph*)]
    \item created intricate pottery and jewelry.
    \item discovered the use of fire.
    \item invented the wheel.
    \item stored food surpluses in good years for distribution in bad years.
    \item established trade routes with distant cultures.
\end{enumerate}
\item Which archaeological method is best for identifying large-scale land modifications?
\begin{enumerate}[label=(\alph*)]
    \item stratigraphic analysis
    \item resistivity survey
    \item metal detectors
    \item ground-penetrating radar
    \item underwater archaeology
\end{enumerate}
\item From which of the following primates do humans descend?
\begin{enumerate}[label=(\alph*)]
    \item bonobos
    \item orangutans
    \item gibbons
    \item chimpanzees
    \item baboons
\end{enumerate}
\end{enumerate}

\section*{True / False (10 marks)}
\textit{Answer True or False}
\begin{enumerate}[label=\arabic*.]
\setcounter{enumi}{5}
\item True or False: Dogs may have provided warmth, as referred to in the Australian Aboriginal expression "three dog night" (an exceptionally cold night), and they would have alerted the camp to the presence of predators or strangers, using their acute hearing to provide an early warning.
\item True or False: Its knowledge and uses spread from China through the Middle East to medieval Europe in the 13 th century, where the first water powered paper mills were built. Because of paper's introduction to the West through the city of Baghdad, it was first called bagdatikos. In the 19 th century, industrial manufacture greatly lowered its cost, enabling mass exchange of information and contributing to significant cultural shifts.
\item True or False: Ashkenazi Jews represent the bulk of modern Jewry, with at least 70\% of Jews worldwide (and up to 90\% prior to World War II and the Holocaust).
\item True or False: The army was still equipped with the Dreyse needle gun of Battle of Königgrätz fame, which was by this time showing the age of its a different answer.
\item True or False: Ashkenazi Jews have won a large number of the Nobel awards. While they make up about 2\% of the U.S. population, a different answer of United States Nobel prize winners in the 20 th century, a quarter of Fields Medal winners, 25\% of ACM Turing Award winners, half the world's chess champions, including 8\% of the top 100 world chess players, and a quarter of Westinghouse Science Talent Search winners have Ashkenazi Jewish ancestry.
\end{enumerate}

\section*{Long Form Response (20 marks)}
\textit{Answer each question with a clear and complete explanation.}
\begin{enumerate}[label=\arabic*.]
\setcounter{enumi}{10}
\item Read the following passage and answer the question: Palermo has at least 2 circuits of City Walls - many pieces of which still survive. The first circuit surrounded the ancient core of the punic City - the so-called Palaeopolis (in the area east of Porta Nuova) and the Neopolis. Via Vittorio Emanuele was the main road east-west through this early walled city. The eastern edge of the walled city was on Via Roma and the ancient port in the vicinity of Piazza Marina. The wall circuit was approximately Porto Nuovo, Corso Alberti, Piazza Peranni, Via Isodoro, Via Candela, Via Venezia, Via Roma, Piazza Paninni, Via Biscottari, Via Del Bastione, Palazzo dei Normanni and back to Porto Nuovo. Question: On what road was the eastern edge of the walled city?
\item Read the following passage and answer the question: Definitions of "Southeast Asia" vary, but most definitions include the area represented by the countries (sovereign states and dependent territories) listed below. All of the states except for East Timor are members of the Association of Southeast Asian Nations (ASEAN). The area, together with part of South Asia, was widely known as the East Indies or simply the Indies until the 20 th century. Christmas Island and the Cocos (Keeling) Islands[citation needed] are considered part of Southeast Asia though they are governed by Australia.[citation needed] Sovereignty issues exist over some territories in the South China Sea. Papua New Guinea has stated that it might join ASEAN, and is currently an observer. Question: What does ASEAN mean?
\end{enumerate}

\vfill
\noindent\textit{End of Paper}
\end{document}